% GENERATED FILE. Edit canonical lessons, then run npm run generate:books.
% canonical-source-hash: fnv1a64:2faa4700e1997330
% canonical-lessons: ZH-C307-fenshou, ZH-C307-lihun, ZH-C307-zaihun, ZH-C307-miyue, ZH-C307-hunsha, ZH-R307-review-302-307

\chapter{Marrying and Parting}
\label{ch:zh-marrying-and-parting}
\hlchaptermodality{307}

\begin{chapteropening}
\textbf{By the end of this chapter:} \emph{I can talk about a honeymoon, breaking up, divorce and marrying again.}

Complete the last lesson of chapter 307: I can talk about a honeymoon, breaking up, divorce and marrying again.
\end{chapteropening}

\section[fēnshǒu]{fēnshǒu — to break up; to end a relationship}

\label{lesson:ZH-C307-fenshou}



\begin{quote}
Before the new one: say the Mandarin for to hold a ceremony, then the Mandarin for wedding sweets for the guests.
\end{quote}

\subsection*{You'll want to know: fēnshǒu}
\emph{fēnshǒu} — \textquotedblleft{}to break up; to end a relationship\textquotedblright{}.

\begin{quote}
\emph{tāmen fēnshǒu le, wǒ hěn nánguò} — They've split up, and I'm sad.
\end{quote}

Its characters come later; for now it is a word for the ear.

\begin{cousinweb}
\emph{fēnshǒu} is built from \emph{fēn} 'to part' and \zh{手} 'a hand', both inherited from Middle Chinese, \emph{pjun} and \emph{syuwX} (Guangyun, in Baxter's transcription), which glosses the first 'to give out' and, quoting an older dictionary, 'to separate': to part hands.
\end{cousinweb}

\begin{grammarlens}[title={What you've built}]
One more word for telling what happened.
\end{grammarlens}

\subsection*{Your turn}
\begin{itemize}
\raggedright
  \item \emph{Say it:} \emph{fēnshǒu}
  \item \emph{Say it:} \emph{fēnshǒu}, once more
  \item \emph{Say it:} \emph{tāmen fēnshǒu le, wǒ hěn nánguò} — They've split up, and I'm sad.
  \item \emph{Recall:} say the Mandarin for to hold a ceremony, then the Mandarin for wedding sweets for the guests, then say \emph{fēnshǒu} again
\end{itemize}

\begin{tcolorbox}[breakable,colback=teal!4,colframe=teal!35!black,title={Before you move on}]
What does fēnshǒu mean? (\textquotedblleft{}To break up; to end a relationship\textquotedblright{}.) Say it once more.
\end{tcolorbox}
\section[líhūn]{líhūn — to divorce}

\label{lesson:ZH-C307-lihun}



\begin{quote}
Before the new one: say the Mandarin for wedding sweets for the guests, then the Mandarin for to end a relationship.
\end{quote}

\subsection*{You'll want to know: líhūn}
\emph{líhūn} — \textquotedblleft{}to divorce\textquotedblright{}.

\begin{quote}
\emph{tāmen jiéhūn shí nián yǐhòu líhūn le} — They divorced after ten years of marriage.
\end{quote}

Its characters come later; for now it is a word for the ear.

\begin{cousinweb}
\emph{líhūn} is built from \zh{离} 'to leave' and \emph{hūn} 'marriage', both inherited from Middle Chinese, \emph{lje} and \emph{xwon} (Guangyun, in Baxter's transcription), which glosses the second 'marriage'.
\end{cousinweb}

\begin{grammarlens}[title={What you've built}]
One more word for telling what happened.
\end{grammarlens}

\subsection*{Your turn}
\begin{itemize}
\raggedright
  \item \emph{Say it:} \emph{líhūn}
  \item \emph{Say it:} \emph{líhūn}, once more
  \item \emph{Say it:} \emph{tāmen jiéhūn shí nián yǐhòu líhūn le} — They divorced after ten years of marriage.
  \item \emph{Recall:} say the Mandarin for wedding sweets for the guests, then the Mandarin for to end a relationship, then say \emph{líhūn} again
\end{itemize}

\begin{tcolorbox}[breakable,colback=teal!4,colframe=teal!35!black,title={Before you move on}]
What does líhūn mean? (\textquotedblleft{}To divorce\textquotedblright{}.) Say it once more.
\end{tcolorbox}
\section[zàihūn]{zàihūn — to remarry}

\label{lesson:ZH-C307-zaihun}



\begin{quote}
Before the new one: say the Mandarin for to end a relationship, then the Mandarin for to divorce.
\end{quote}

\subsection*{You'll want to know: zàihūn}
\emph{zàihūn} — \textquotedblleft{}to remarry\textquotedblright{}.

\begin{quote}
\emph{tā líhūn yǐhòu, hěn kuài jiù zàihūn le} — After his divorce he soon married again.
\end{quote}

Its characters come later; for now it is a word for the ear.

\begin{cousinweb}
\emph{zàihūn} is built from \zh{再} 'again' and \emph{hūn} 'marriage', both inherited from Middle Chinese, \emph{tsojH} and \emph{xwon} (Guangyun, in Baxter's transcription), which glosses the first 'twice'.
\end{cousinweb}

\begin{grammarlens}[title={What you've built}]
One more word for telling what happened.
\end{grammarlens}

\subsection*{Your turn}
\begin{itemize}
\raggedright
  \item \emph{Say it:} \emph{zàihūn}
  \item \emph{Say it:} \emph{zàihūn}, once more
  \item \emph{Say it:} \emph{tā líhūn yǐhòu, hěn kuài jiù zàihūn le} — After his divorce he soon married again.
  \item \emph{Recall:} say the Mandarin for to end a relationship, then the Mandarin for to divorce, then say \emph{zàihūn} again
\end{itemize}

\begin{tcolorbox}[breakable,colback=teal!4,colframe=teal!35!black,title={Before you move on}]
What does zàihūn mean? (\textquotedblleft{}To remarry\textquotedblright{}.) Say it once more.
\end{tcolorbox}
\section[mìyuè]{mìyuè — honeymoon}

\label{lesson:ZH-C307-miyue}



\begin{quote}
Before the new one: say the Mandarin for to divorce, then the Mandarin for to remarry.
\end{quote}

\subsection*{You'll want to know: mìyuè}
\emph{mìyuè} — \textquotedblleft{}honeymoon\textquotedblright{}.

\begin{quote}
\emph{mìyuè de shíhou, wǒmen qù le hǎibiān} — On our honeymoon we went to the seaside.
\end{quote}

Its characters come later; for now it is a word for the ear.

\begin{cousinweb}
\emph{mìyuè} is built from \emph{mì} 'honey' and \zh{月} 'a month', both inherited from Middle Chinese, \emph{mjit} and \emph{ngjwot} (Guangyun, in Baxter's transcription), which glosses the first as the food that bees make: the sweet month after a wedding.
\end{cousinweb}

\begin{grammarlens}[title={What you've built}]
One more word for telling what happened.
\end{grammarlens}

\subsection*{Your turn}
\begin{itemize}
\raggedright
  \item \emph{Say it:} \emph{mìyuè}
  \item \emph{Say it:} \emph{mìyuè}, once more
  \item \emph{Say it:} \emph{mìyuè de shíhou, wǒmen qù le hǎibiān} — On our honeymoon we went to the seaside.
  \item \emph{Recall:} say the Mandarin for to divorce, then the Mandarin for to remarry, then say \emph{mìyuè} again
\end{itemize}

\begin{tcolorbox}[breakable,colback=teal!4,colframe=teal!35!black,title={Before you move on}]
What does mìyuè mean? (\textquotedblleft{}Honeymoon\textquotedblright{}.) Say it once more.
\end{tcolorbox}
\section[hūnshā]{hūnshā — wedding dress}

\label{lesson:ZH-C307-hunsha}



\begin{quote}
Before the new one: say the Mandarin for to remarry, then the Mandarin for a honeymoon.
\end{quote}

\subsection*{You'll want to know: hūnshā}
\emph{hūnshā} — \textquotedblleft{}wedding dress\textquotedblright{}.

\begin{quote}
\emph{xīnniáng de hūnshā hěn piàoliang} — The bride's dress is beautiful.
\end{quote}

Its characters come later; for now it is a word for the ear.

\begin{cousinweb}
\emph{hūnshā} is built from \emph{hūn} 'marriage' and \emph{shā} 'gauze', both inherited from Middle Chinese, \emph{xwon} and \emph{srae} (Guangyun, in Baxter's transcription), which glosses the second 'a kind of silk'.
\end{cousinweb}

\begin{grammarlens}[title={What you've built}]
One more word for telling what happened.
\end{grammarlens}

\subsection*{Your turn}
\begin{itemize}
\raggedright
  \item \emph{Say it:} \emph{hūnshā}
  \item \emph{Say it:} \emph{hūnshā}, once more
  \item \emph{Say it:} \emph{xīnniáng de hūnshā hěn piàoliang} — The bride's dress is beautiful.
  \item \emph{Recall:} say the Mandarin for to remarry, then the Mandarin for a honeymoon, then say \emph{hūnshā} again
\end{itemize}

\begin{tcolorbox}[breakable,colback=teal!4,colframe=teal!35!black,title={Before you move on}]
What does hūnshā mean? (\textquotedblleft{}Wedding dress\textquotedblright{}.) Say it once more.
\end{tcolorbox}
\section[(dialogue)]{Looking back: from A Baby on the Way to Marrying and Parting}

\label{lesson:ZH-R307-review-302-307}



\begin{quote}
Say the Mandarin for to be expecting a baby and to stay alive.
\end{quote}

\subsection*{Your turn}
\begin{itemize}
\raggedright
  \item \emph{Say it:} to be expecting a baby, to stay alive, a newborn, to bring back to life, a pair of twins
  \item \emph{Say it:} to name a baby, a childhood pet name, a kindergarten, to start the school term, a new student at school
  \item \emph{Say it:} to start school as a child, a diary, to report in on arrival, school tuition fees, to commit to memory
  \item \emph{Say it:} to go out with someone, to propose marriage, to get engaged, a happy occasion, the bridegroom
  \item \emph{Say it:} the bride, the whole story of what happened, the bride and groom, to hold a ceremony, wedding sweets for the guests
  \item \emph{Say it:} to end a relationship, to divorce, to remarry, a honeymoon, a white wedding dress
  \item \emph{Read it:} \zh{活} \zh{记} \zh{经} — each one aloud, then a word that holds it
\end{itemize}

\begin{tcolorbox}[breakable,colback=teal!4,colframe=teal!35!black,title={Before you move on}]
To be expecting a baby? (\emph{huáiyùn}.) To stay alive? (\textbf{\zh{活}}, \emph{huó}.) A newborn? (\emph{yīng'ér}.) To bring back to life? (\textbf{\zh{救活}}, \emph{jiùhuó}.) A pair of twins? (\emph{shuāngbāotāi}.) To name a baby? (\textbf{\zh{起名}}, \emph{qǐmíng}.) A childhood pet name? (\textbf{\zh{小名}}, \emph{xiǎomíng}.) A kindergarten? (\emph{yòu'éryuán}.) To start the school term? (\textbf{\zh{开学}}, \emph{kāixué}.) A new student at school? (\textbf{\zh{新生}}, \emph{xīnshēng}.) To start school as a child? (\emph{rùxué}.) A diary? (\textbf{\zh{日记}}, \emph{rìjì}.) To report in on arrival? (\textbf{\zh{报到}}, \emph{bàodào}.) School tuition fees? (\emph{xuéfèi}.) To commit to memory? (\emph{jìzhù}.) To go out with someone? (\emph{tán liàn'ài}.) To propose marriage? (\emph{qiúhūn}.) To get engaged? (\emph{dìnghūn}.) A happy occasion? (\textbf{\zh{喜事}}, \emph{xǐshì}.) The bridegroom? (\emph{xīnláng}.) The bride? (\emph{xīnniáng}.) The whole story of what happened? (\textbf{\zh{经过}}, \emph{jīngguò}.) The bride and groom? (\textbf{\zh{新人}}, \emph{xīnrén}.) To hold a ceremony? (\emph{jǔxíng}.) Wedding sweets for the guests? (\emph{xǐtáng}.) To end a relationship? (\emph{fēnshǒu}.) To divorce? (\emph{líhūn}.) To remarry? (\emph{zàihūn}.) A honeymoon? (\emph{mìyuè}.) A white wedding dress? (\emph{hūnshā}.)
\end{tcolorbox}
